\selectlanguage{english}

\section{Hodge theory}

To study the topological properties of a smooth manifold through its cohomology groups, it is convenient to define an additional operation on forms: the Hodge dual.

\subsection{Hodge operator}

\begin{definition}{Hodge dual}{}
  Given an $ n $-dimensional oriented metric manifold $ (\mathcal{M},g) $, the \bcdef{Hodge dual} is defined as the map $ \star : \lm{p} \rightarrow \lm{n-p} : \omega \mapsto \star\, \omega $ such that:
  \begin{equation}
    (\star\, \omega)_{\mu_1 \dots \mu_{n-p}} \defeq \frac{1}{p!} \sqg\, \epsilon_{\nu_1 \dots \nu_p \mu_1 \dots \mu_{n-p}} \omega^{\nu_1 \dots \nu_p}
    \label{eq:hodge-dual}
  \end{equation}
\end{definition}

The factorials need to be treated carefully:
\begin{equation}
  \star\, \omega \equiv \frac{1}{(n-p)!} (\star\, \omega)_{\mu_1 \dots \mu_{n-p}} \dd x^{\mu_1} \wedge \dots \wedge \dd x^{\mu_{n-p}}
\end{equation}

In this section, the orientedness and $ n $-dimensionality of the manifold $ \mathcal{M} $ are implied.

\begin{proposition}{Invariance of the Hodge dual}{}
  The Hodge dual is basis-independent and conformally-invariant.
\end{proposition}

It is useful to state a lemma for future calculations.

\begin{lemma}{Generalized Kronecker delta}{}
  $ v^{\mu_1 \dots \mu_p \rho_1 \dots \rho_{n-p}} v_{\nu_1 \dots \nu_p \rho_1 \dots \rho_{n-p}} = \sigma p! (n-p)! \delta^{\mu_1}_{[\nu_1} \dots \delta^{\mu_p}_{\nu_p]} $.
\end{lemma}

\begin{proposition}{Nilpotence}{hodge-hodge}
  $ \star \left( \star\, \omega \right) = \sigma \left( -1 \right)^{p (n-p)} \omega $.
\end{proposition}

The Hodge dual defines an inner product on each $ \lm{p} $:
\begin{equation}
  \braket{\omega, \eta} \defeq \int_{\mathcal{M}} \omega \wedge \star\, \eta
\end{equation}
This allows to define operators and their adjoints on form spaces.

\begin{proposition}{Adjoint exterior derivative}{}
  Given a metric manifold $ (\mathcal{M},g) $ and two forms $ \omega \in \lm{p}, \alpha \in \lm{p-1} $, then:
  \begin{equation}
    \braket{\dd \alpha,\omega} = \braket{\alpha,\dd^{\dagger}\omega}
  \end{equation}
  where the adjoint of the exterior derivative $ \dd^{\dagger} : \lm{p} \rightarrow \lm{p-1} $ is defined as:
  \begin{equation}
    \dd^{\dagger} \defeq \sigma \left( -1 \right)^{np + n -1} \star \dd \, \star
  \end{equation}
\end{proposition}

\begin{proofbox}
  \begin{proof}
    To simplify the proof, consider a closed manifold; then, from Stokes' theorem and \eref{eq:ext-der-prod-rule}:
    \begin{equation*}
      0 = \int_{\mathcal{M}} \dd ( \alpha \wedge \star\, \omega ) = \braket{\dd\alpha, \omega} + \int_{\mathcal{M}} \left( -1 \right)^{p-1} \alpha \wedge \dd \star \omega
    \end{equation*}
    The second term is proportional to $ \braket{\alpha, \star\, \dd \star \omega} $: to determine the relative sign, note that $ \dd \star \omega \in \lm{n-p+1} $, thus, from \pref{prop:hodge-hodge}, $ \star \star \dd \star \omega = \sigma \left( -1 \right)^{(n-p+1)(p-1)} \dd \star \omega $. In conclusion:
    \begin{multline*}
      \braket{\alpha, \star\, \dd \star \omega} = \sigma \left( -1 \right)^{(n-p)(p-1)} \int_{\mathcal{M}} \left( -1 \right)^{p-1} \alpha \wedge \dd \star \omega \\
      \quad \implies \quad
      \braket{d\alpha, \omega} = \sigma \left( -1 \right)^{(n-p)(p-1) + 1} \braket{\alpha, \star\, \dd \star \omega}
    \end{multline*}
    Noting that $ \left( -1 \right)^{(n-p)(p-1) + 1} = \left( -1 \right)^{np + n - 1} $, as in general $ \left( -1 \right)^{-n} = \left( -1 \right)^n $ and $ \left( -1 \right)^{-p^2 + p + 1} = \left( -1 \right)^{-1} $ due to $ p (p-1) $ being always even, concludes the proof.
  \end{proof}
\end{proofbox}

\begin{definition}{Laplace--Beltrami operator}{}
  Given a metric manifold $ (\mathcal{M},g) $, the \bcdef{Laplace--Beltrami operator} $ \lap : \lm{p} \rightarrow \lm{p} $ is defined as the operator:
  \begin{equation}
    \lap \defeq ( \dd + \dd^{\dagger} )^2
  \end{equation}
\end{definition}

Clearly, since $ \dd^2 = {\dd\dg}^2 = 0 $:
\begin{equation}
  \lap = \dd\dd^{\dagger} + \dd^{\dagger}\dd \equiv \{\dd,\dd^{\dagger}\}
\end{equation}

It is possible to compute an explicit expression for the action of the Laplace--Beltrami operator on functions (usually called the ``Laplacian").

\begin{lemma}{}{lap-calc}
  Given $ f \in \cm $, then $ \dd^{\dagger} f = 0 $.
\end{lemma}

\begin{proofbox}
  \begin{proof}
    Trivial noting that $ \star\, f $ is a top-form.
  \end{proof}
\end{proofbox}

\begin{proposition}{}{}
  Given $ f \in \cm $, then:
  \begin{equation}
    \lap f = - \frac{\sigma}{\sqg} \pa_{\nu} \left( \sqg\, g^{\mu \nu} \pa_{\mu} f \right)
  \end{equation}
\end{proposition}

\begin{proofbox}
  \begin{proof}
    Via direct calculation, using \lref{lemma:lap-calc}:
    \begin{equation*}
      \begin{split}
        \lap f
        &= \sigma \left( -1 \right)^{2n - 1} \star \dd \star \left( \pa_{\mu} f \dd x^{\mu} \right) = - \sigma \star \dd \left( \pa_{\mu} f \star \dd x^{\mu} \right) \\
        &= - \frac{\sigma}{(n-1)!} \star \dd \left( \pa_{\mu} f g^{\mu \nu} \sqg\, \epsilon_{\nu \rho_1 \dots \rho_{n-1}} \dd x^{\rho_1} \wedge \dots \wedge \dd x^{\rho_{n-1}} \right) \\
        &= - \frac{\sigma}{(n-1)!} \star \pa_{\alpha} \left( \sqg\, g^{\mu \nu} \pa_{\mu} f \right) \epsilon_{\nu \rho_1 \dots \rho_{n-1}} \dd x^{\alpha} \wedge \dd x^{\rho_1} \wedge \dots \wedge \dd x^{\rho_{n-1}} \\
        &= - \sigma \star \pa_{\nu} \left( \sqg\, g^{\mu \nu} \pa_{\mu} f \right) \dd x^1 \wedge \dots \dd x^n = - \frac{\sigma}{\sqg} \pa_{\nu} \left( \sqg\, g^{\mu \nu} \pa_{\mu} f \right)
      \end{split}
    \end{equation*}
    which is the thesis.
  \end{proof}
\end{proofbox}

\subsection{de Rham cohomology}
\label{ssec:derham}

While exact $ \implies $ closed, the converse is not true, in general: it depends on the topological properties of the manifold.

\begin{lemma}{Poincaré's lemma}{}
  If $ \mathcal{M} $ is simply connected, then $ \omega\in\lm{p} $ closed $ \implies \omega $ exact.
\end{lemma}

In general, it is always possible to choose a simply connected neighbourhood of a point $ p \in \mathcal{M} $, in which every closed form is exact, but that may not always be possible globally.

It is convenient to set the notation $ d_p \equiv d : \lm{p} \rightarrow \lm{p+1} $, so that the set of all closed $ p $-forms on $ \mathcal{M} $ is denoted by $ Z^p(\mathcal{M}) \defeq \ker d_p $, while the set of all exact $ p $-forms by $ B^p(\mathcal{M}) \defeq \ran d_{p-1} $.

\begin{definition}{Equivalent forms}{}
  Two closed $ p $-forms $ \omega,\omega' \in Z^P(\mathcal{M}) $ are said to be \bcdef{equivalent} if $ \omega = \omega' + \eta $ for some $ \eta \in B^p(\mathcal{M}) $.
\end{definition}

\begin{definition}{de Rahm cohomology}{}
  The \bcdef{$ p^{\text{th}} $ de Rham cohomolgy group} of a manifold $ \mathcal{M} $ is defined to be:
  \begin{equation}
    H^p(\mathcal{M}) \defeq Z^p(\mathcal{M}) / B^p(\mathcal{M})
  \end{equation}
  and its dimension is the $ p^\text{th} $ \bcdef{Betti number} of $ \mathcal{M} $:
  \begin{equation}
    B_p \defeq \dim_{\R} H^p(\mathcal{M})
  \end{equation}
\end{definition}

\begin{theorem}{Betti numbers}{}
  Given a smooth manifold, its Betti numbers are always finite.
\end{theorem}

$ B_0 = 1 $ for any connected manifold: there exist constant functions, which are manifestly closed and not exact, due to the non-existence of ``$ (-1) $-forms". Higher Betti numbers are non-zero only if the manifold has some non-trivial topology.

\begin{definition}{Eulerv's character}{}
  The \bcdef{Euler's character} of a manifold $ \mathcal{M} $ is defined as:
  \begin{equation}
    \chi(\mathcal{M}) \defeq \sum_{p \in \N_0} (-1)^p B_p
  \end{equation}
\end{definition}

\begin{example}{Spheres and tori}{}
  The $ n $-sphere $ \mathbb{S}^n $ has only non-vanishing $ B_0 = B_n = 1 $, thus $ \chi(\mathbb{S}^n) = 1 + (-1)^n $, while the $ n $-torus $ \mathbb{T}^n $ has $ B_p = \binom{n}{p} $, hence $ \chi(\mathbb{T}^n) = 0 $.
\end{example}

There is an additional characterization of the de Rham cohomology through the Laplace--Beltrami operator.

\begin{definition}{Harmonic $ p $-forms}{}
  Given $ \omega \in \lm{p} $, it is said to be \bcdef{harmonic} if $ \lap \omega = 0 $.
  The space of harmonic $ p $-forms on $ (\mathcal{M},g) $ is denoted as $ \hrm{p} $.
\end{definition}

\begin{proposition}{}{harm-closed}
  A harmonic form is both \bcprop{closed} and \bcprop{co-closed}.
\end{proposition}

\begin{proofbox}
  \begin{proof}
    $ 0 = \braket{\omega, \lap \omega} = \braket{\dd\omega, \dd\omega} + \braket{\dd^{\dagger} \omega, \dd^{\dagger}{\omega}} $, thus $ \dd\omega = 0 $ and $ \dd^{\dagger} \omega = 0 $, for the inner product is positive-defined.
  \end{proof}
\end{proofbox}

\begin{theorem}{Harmonic decomposition of forms}{form-decomp}
  Given a compact Riemannian manifold $ (\mathcal{M},g) $, any $ \omega \in \lm{p} $ can be uniquely decomposed as $ \omega = \dd \alpha + \dd^{\dagger}\beta + \gamma $, with $ \alpha \in \lm{p-1} $, $ \beta \in \lm{p+1} $ and $ \gamma \in \hrm{p} $.
\end{theorem}

\begin{theorem}{Hodge's theorem}{}
  Given a compact Riemannian manifold $ (\mathcal{M},g) $, then:
  \begin{equation}
    \hrm{p} \cong H^p(\mathcal{M})
  \end{equation}
\end{theorem}

\begin{proofbox}
  \begin{proof}
    From \pref{prop:harm-closed} $ \hrm{p} \subset Z^p(\mathcal{M}) $, but the uniqueness of decomposition in \tref{th:form-decomp} implies $ \forall \gamma \in \hrm{p} \,\exists \eta_{\gamma} \in \lm{p-1} : \gamma \neq \dd\eta_{\gamma} $, thus $ \hrm{p} \subset H^p(\mathcal{M}) $.

    WTS that any equivalence class $ [\omega] \in H^p(\mathcal{M}) $ can be represented by a harmonic form. By \tref{th:form-decomp} $ \omega = \dd\alpha + \dd^{\dagger}\beta + \gamma $, but $ \omega \in H^p(\mathcal{M}) $ implies $ \dd\omega = 0 $ by definition, so:
    \begin{equation*}
      0 = \braket{\dd\omega, \beta} = \braket{\omega, \dd^{\dagger} \beta} = \braket{\dd\alpha + \dd^{\dagger}\beta + \gamma, \dd^{\dagger}\beta} = \braket{\dd^{\dagger} \beta, \dd^{\dagger} \beta}
    \end{equation*}
    The inner product is positive-definite, thus $ \dd^{\dagger} \beta = 0 $, hence $ \omega = \gamma + \dd\alpha $. By definition $ H^p(\mathcal{M}) \defeq Z^p(\mathcal{M}) / B^p(\mathcal{M}) $, so $ [\omega] = \gamma $.
  \end{proof}
\end{proofbox}

Betti numbers can then be computed as $ B_p = \dim_{\R} \hrm{p} $.

\newpage
\section{Cartan formalism}

Although typically calculations are carried on a coordinate basis $ \{e_\mu\} = \{\pa_\mu\} $, it is possible to choose bases without such an interpretation. For example, a linear combination of a coordinate basis is not in general a coordinate basis itself:
\begin{equation}
  e_a = \tensor{e}{_a^\mu} \pa_\mu
\end{equation}
with $ a = 1, \dots, n $ like $ \mu = 1, \dots, n $.
A particularly useful non-coordinate basis is one such that:
\begin{equation}
  g(e_a, e_b) = g_{\mu \nu} \tensor{e}{_a^\mu} \tensor{e}{_b^\nu} = \eta_{ab}
  \label{eq:vielbeins-def}
\end{equation}
The components $ \tensor{e}{_a^\mu} $ are called \bctxt{vielbeins} or \textit{tetrads}. In this non-coordinate system, the manifold looks flat (or, at least, its patch covered by the given chart). In the following computations, Greek indices are raised/lowered by the metric $ g_{\mu \nu} $, while Latin indices by the metric $ \eta_{ab} $.

The vielbeins are not unique: given $ \{\tensor{e}{_a^\mu}\} $, it is always possible to find a new vielbeins like:
\begin{equation}
  \tensor{\tilde{e}}{_a^\mu} = \tensor{e}{_b^\mu} \tensor{(\Lambda^{-1})}{^b_a}
  \label{eq:vielbeins-trans}
\end{equation}
The transformation matrix must satisfy the condition imposed by \eref{eq:vielbeins-def}, i.e.:
\begin{equation}
  \tensor{\Lambda}{^c_a} \tensor{\Lambda}{^d_b} \eta_{cd} = \eta_{ab}
\end{equation}
These are \bctxt{local Lorentz transformations}, because this condition defines an $ \lorgn $ transformation\footnotemark, but $ \Lambda $ is now allowed to vary over the manifold. The dual basis of 1-forms $ \{\theta^a\} $ is defined by $ \theta^a(e_b) = \delta^a_b $. The relation to the coordinate basis is:
\begin{equation}
  \theta^a = \tensor{e}{^a_\mu} \dd x^\mu
\end{equation}
where the coefficients satisfy $ \tensor{e}{^a_\mu} \tensor{e}{_b^\mu} = \delta^a_b $ and $ \tensor{e}{^a_\mu} \tensor{e}{_a^\nu} = \delta_\mu^\nu $.
The metric is a tensor, therefore $ g = g_{\mu \nu} \dd x^\mu \otimes \dd x^\nu = \eta_{ab} \theta^a \otimes \theta^b $, thus it is related to the vielbeins by:
\begin{equation}
  g_{\mu \nu} = \tensor{e}{^a_\mu} \tensor{e}{^b_\nu} \eta_{ab}
\end{equation}

\footnotetext{To be precise, this condition defines a $ \On{1,n} $ transformation on $ \R^{1,n} $, but this orthogonal group has four disconnected components: the $ \lorgn $ component is the one defined by $ \det \Lambda = 1 $ and $ \tensor{\Lambda}{^0_0} \ge 1 $, while the other three components can be obtained by composition of $ \lorgn $ with the possible combinations of parity reversal and time reversal. For a more detailed account on the Lorentz group, see \S 1.2 of \cite{leonardo}.}

\subsection{Connection 1-form}

On a vielbein basis $ \{e_a\} $, connection components are computed in the usual way:
\begin{equation}
  \na_{e_c} e_b \equiv \omega^a_{cb} e_a
  \label{eq:connection-vielbeins}
\end{equation}
where the $ n^3 $ functions $ \omega^a_{cb} $ are just the connection coefficients $ \Gamma^a_{cb} $ on the vielbein basis.

\begin{definition}{Connection 1-form}{}
  On a metric manifold with vielbeins, the \bcdef{connection 1-form} is defined as:
  \begin{equation}
    \tensor{\omega}{^a_b} \defeq \omega^a_{cb} \theta^c
    \label{eq:connection-1-form}
  \end{equation}
\end{definition}

Note that these are really $ n^2 $ 1-forms, according to values of $ a,b = 1, \dots, n $. This is also known as the \textit{spin connection}, due to its relationship to spinors in curved spacetime.
Then, the connection in the vielbein basis is easily defined:
\begin{equation*}
  \na_X Y = X(Y^a) e_a + Y^a X^b \na_{e_b} e_a = X^b [e_b(Y^a) + \omega^a_{bc} Y^c] e_a
\end{equation*}
Hence, setting the notation $ \na_b Y^a \equiv (\na_{e_b} Y)^a $, it is clear that the connection takes the form:
\begin{equation}
  \na_b Y^a = e_b(Y^a) + \omega^a_{bc} Y^c
  \label{eq:vielbein-cov-der}
\end{equation}
In general, $ e_b(Y^a) = \tensor{e}{_b^\mu} \pa_\mu Y^a $. Using \eref{eq:cov-der-form-aux}, the covariant derivative of forms is found:
\begin{equation}
  \na_b \zeta_a = e_b(\zeta_a) - \omega^c_{ba} \zeta_c
\end{equation}
Finally, for tensors of general rank:
\begin{equation}
  \na_c \tensor{T}{^{a_1 \dots a_r}_{b_1 \dots b_s}} = e_c(\tensor{T}{^{a_1 \dots a_r}_{b_1 \dots b_s}}) + \sum_{j = 1}^r \omega^{a_j}_{c d_j} \tensor{T}{^{a_1 \dots d_j \dots a_n}_{b_1 \dots b_s}} - \sum_{j = 1}^s \omega^{d_j}_{c b_j} \tensor{T}{^{a_1 \dots a_n}_{b_1 \dots d_j \dots b_n}}
\end{equation}

\begin{lemma}{Transformation law}{}
  Given a local Lorentzian transformation $ \Lambda $ acting as in \eref{eq:vielbeins-trans}:
  \begin{equation}
    \tensor{\tilde{\omega}}{^a_b} = \tensor{\Lambda}{^a_c} \tensor{\omega}{^c_d} \tensor{(\Lambda^{-1})}{^d_b} + \tensor{\Lambda}{^a_c} \dd \tensor{(\Lambda^{-1})}{^c_b}
    \label{eq:con-1-form-trans}
  \end{equation}
\end{lemma}

\begin{proofbox}
  \begin{proof}
    For ease of notation, rename the transformation in \eref{eq:vielbeins-trans} as $ \Lambda^{-1} \equiv \mt{L} $. Then, from the orthogonality condition $ \theta^a (e_b) = \delta^a_b $, it is clear that:
    \begin{equation*}
      \tilde{e}_a = \tensor{\mt{L}}{^b_a} e_b
      \qquad \qquad
      \tilde{\theta}^a = \tensor{(\mt{L}^{-1})}{^a_b} \theta^b
    \end{equation*}
    Then, the connection coefficients can be computed as:
    \begin{equation*}
      \begin{split}
        \tilde{\omega}^a_{cb} \tilde{e}_a
        & = \na_{\tilde{e}_c} \tilde{e}_b = \tensor{\mt{L}}{^d_c} \na_{e_d} (\tensor{\mt{L}}{^f_b} e_f) = \tensor{\mt{L}}{^d_c} \left[ \na_{e_d} \tensor{\mt{L}}{^f_b} e_f + \tensor{\mt{L}}{^f_b} \na_{e_d} e_f \right] \\
        & = \tensor{\mt{L}}{^d_c} \left[ \na_{e_d} \tensor{\mt{L}}{^f_b} e_f + \tensor{\mt{L}}{^f_b} \omega^g_{df} e_g \right] = \tensor{\mt{L}}{^d_c} \left[ \na_{e_d} \tensor{\mt{L}}{^f_b} + \tensor{\mt{L}}{^g_b} \omega^f_{dg} \right] \tensor{(\mt{L}^{-1})}{^a_f} \tilde{e}_a
      \end{split}
    \end{equation*}
    Finally, the transformed connection 1-form is:
    \begin{equation*}
      \begin{split}
        \tensor{\tilde{\omega}}{^a_b}
        & = \tilde{\omega}^a_{cb} \tilde{\theta^c} = \tensor{\mt{L}}{^d_c} \left[ \na_{e_d} \tensor{\mt{L}}{^f_b} + \tensor{\mt{L}}{^g_b} \omega^f_{dg} \right] \tensor{(\mt{L}^{-1})}{^a_f} \tensor{(\mt{L}^{-1})}{^c_h} \theta^h \\
        & = \left[ \na_{e_d} \tensor{\mt{L}}{^f_b} + \tensor{\mt{L}}{^g_b} \omega^f_{dg} \right] \tensor{(\mt{L}^{-1})}{^a_f} \theta^d = \tensor{(\mt{L}^{-1})}{^a_c} \left[ \dd x^\mu \na_\mu \tensor{\mt{L}}{^c_b} + \tensor{\omega}{^c_d} \tensor{\mt{L}}{^d_b} \right]
      \end{split}
    \end{equation*}
    Recalling that $ \dd f = \dd x^\mu \na_\mu f $, with $ f \in \cm $, and that $ \Lambda = \mt{L}^{-1} $ yields the thesis.
  \end{proof}
\end{proofbox}

The second term reflects the second term in \pref{prop:gamma-non-tens}, which involves the derivative of the coordinate transformation.

As for the connection coefficients, it is possible to compute the torsion tensor on a vielbein basis as $ \tensor{T}{^a_{bc}} = \tors(\theta^a , e_b , e_c) $, which retains the antisimmetry $ \tensor{T}{^a_{bc}} = - \tensor{T}{^a_{cb}} $.

\begin{definition}{Torsion 2-form}{}
  On a metric manifold with vielbeins, the \bcdef{torsion 2-form} is defined as:
  \begin{equation}
    T^a \defeq \frac{1}{2} \tensor{T}{^a_{bc}} \theta^b \wedge \theta^c
  \end{equation}
\end{definition}

This is a vector-valued 2-form, thus $ T^a $ has $ n $ independent 2-forms. Like \eref{eq:torsion-comp}, the torsion 2-form is linked to the connection 1-form by the first Cartan structure equation.

\begin{theorem}[before upper = {\tcbtitle}]{First Cartan structure equation}{}
  \begin{equation}
    T^a = \dd\theta^a + \tensor{\omega}{^a_b} \wedge \theta^b
  \end{equation}
\end{theorem}

\begin{proofbox}
  \begin{proof}
    From \eref{eq:connection-vielbeins} $ \omega^a_{cb} = \tensor{e}{^a_\rho} \tensor{e}{_c^\mu} \na_\mu \tensor{e}{_b^\rho} $, thus inserting \eref{eq:covariant-derivative-vector} and \eref{eq:connection-1-form}:
    \begin{equation*}
      \begin{split}
        \tensor{\omega}{^a_b} \wedge \theta^b
        & = \omega^a_{cb} \left( \tensor{e}{^c_\mu} \dd x^\mu \right) \wedge \left( \tensor{e}{^b_\nu} \dd x^\nu \right) = \tensor{e}{^a_\rho} \tensor{e}{_c^\lambda} \left( \pa_\lambda \tensor{e}{_b^\rho} + \tensor{e}{_b^\sigma} \Gamma^\rho_{\lambda \sigma} \right) \left( \tensor{e}{^c_\mu} \dd x^\mu \right) \wedge \left( \tensor{e}{^b_\nu} \dd x^\nu \right) \\
        & = \tensor{e}{^a_\rho} \tensor{e}{_c^\lambda} \tensor{e}{^c_\mu} \tensor{e}{^b_\nu} \left( \pa_\lambda \tensor{e}{_b^\rho} + \tensor{e}{_b^\sigma} \Gamma^\rho_{\lambda \sigma} \right) \dd x^\mu \wedge \dd x^\nu = \tensor{e}{^a_\rho} \tensor{e}{^b_\nu} \left( \pa_\mu \tensor{e}{_b^\rho} + \tensor{e}{_b^\sigma} \Gamma^\rho_{\mu \sigma} \right) \dd x^\mu \wedge \dd x^\nu \\
      \end{split}
    \end{equation*}
    Now, $ \tensor{e}{^b_\nu} \tensor{e}{_b^\rho} = \delta_\nu^\rho $ implies $ \tensor{e}{^b_\nu} \pa_\mu \tensor{e}{_b^\rho} = - \tensor{e}{_b^\rho} \pa_\mu \tensor{e}{^n_\nu} $, hence:
    \begin{equation*}
      \begin{split}
        \tensor{\omega}{^a_b} \wedge \theta^b
        & = \left[ - \tensor{e}{^a_\rho} \tensor{e}{_b^\rho} \pa_\mu \tensor{e}{^b_\nu} + \tensor{e}{^a_\rho} \tensor{e}{^b_\nu} \tensor{e}{_b^\sigma} \Gamma^\rho_{\mu \sigma} \right] \dd x^\mu \wedge \dd x^\nu \\
        & = - \pa_\mu \tensor{e}{^a_\nu} \dd x^\mu \wedge \dd x^\nu + \tensor{e}{^a_\rho} \Gamma^\rho_{\mu \nu} \dd x^\mu \wedge \dd x^\nu = -\dd\theta^a + \tensor{e}{^a_\rho} \Gamma^\rho_{[\mu \nu]} \dd x^\mu \wedge \dd x^\nu
      \end{split}
    \end{equation*}
    Now, from the multilinearity of the torsion tensor:
    \begin{equation*}
      \begin{split}
        T^a
        & = \frac{1}{2} \tensor{T}{^a_{bc}} \theta^b \wedge \theta^c = \frac{1}{2} \tensor{e}{^a_\rho} \tensor{e}{_b^\alpha} \tensor{e}{_c^\beta} \tensor{e}{^b_\mu} \tensor{e}{^c_\nu} \tensor{T}{^\rho_{\beta \gamma}} \dd x^\mu \wedge \dd x^\nu = \frac{1}{2} \tensor{e}{^a_\rho} \tensor{T}{^\rho_{\mu \nu}} \dd x^\mu \wedge \dd x^\nu
      \end{split}
    \end{equation*}
    Writing $ \tensor{T}{^\rho_{\mu \nu}} = 2 \Gamma^\rho_{[\mu \nu]} $ completes the proof.
  \end{proof}
\end{proofbox}

This construction can be generalized to arbitrary vector-valued 1-forms, defining the so-called \bctxt{covariant exterior derivative}:
\begin{equation}
  \dcov \zeta^a \defeq \dd \zeta^a + \tensor{\omega}{^a_b} \wedge \zeta^b
\end{equation}
This is a map $ \dcov : \lm{p} \ra \lm{p+1} $, like the exterior derivative, but it has an important invariance property.

\begin{lemma}{Lorentz invariance}{form-lor-inv}
  The covariant exterior derivative transformations under local Lorentz as:
  \begin{equation}
    \zeta^a \mapsto \tensor{\Lambda}{^a_b} \zeta^b
    \quad \implies \quad
    \dcov \zeta^a \mapsto \tensor{\Lambda}{^a_b} \dcov \zeta^b
  \end{equation}
\end{lemma}

\begin{proofbox}
  \begin{proof}
    Denoting transformed quantities with a tilde and recalling \eref{eq:con-1-form-trans}:
    \begin{equation*}
      \begin{split}
        \tilde{\dcov} \tilde{\zeta^a}
        & = \dd \tilde{\zeta^a} + \tensor{\tilde{\omega}}{^a_b} \wedge \tilde{\zeta}^a \\
        & = \dd \tensor{\Lambda}{^a_b} \wedge \zeta^b + \tensor{\Lambda}{^a_b} \dd \zeta^b + \tensor{\Lambda}{^a_c} \tensor{(\Lambda^{-1})}{^d_b} \tensor{\Lambda}{^b_f} \tensor{\omega}{^c_d} \wedge \zeta^f + \tensor{\Lambda}{^a_c} \tensor{\Lambda}{^b_d} \dd \tensor{(\Lambda^{-1})}{^c_b} \wedge \zeta^d \\
        & = \dd \tensor{\Lambda}{^a_b} \wedge \zeta^b + \tensor{\Lambda}{^a_b} \dd \zeta^b + \tensor{\Lambda}{^a_c} \tensor{\omega}{^c_d} \wedge \zeta^d - \tensor{\Lambda}{^a_c} \tensor{(\Lambda^{-1})}{^c_b} \dd \tensor{\Lambda}{^b_d} \wedge \zeta^d = \tensor{\Lambda}{^a_b} \left[ \dd \zeta^b + \tensor{\omega}{^b_c} \zeta^c \right]
      \end{split}
    \end{equation*}
    where $ \tensor{(\Lambda^{-1})}{^c_b} \tensor{\Lambda}{^b_d} = \delta^c_d \implies \tensor{\Lambda}{^b_d} \dd \tensor{(\Lambda^{-1})}{^c_b} = - \tensor{(\Lambda^{-1})}{^c_b} \dd \tensor{\Lambda}{^b_d} $ has been used.
  \end{proof}
\end{proofbox}

The exterior covariant derivative can be used to extend the covariant derivative to vector-valued 1-forms (using \eref{eq:vielbein-cov-der}):
\begin{equation}
  \dcov \zeta^a = \pa_\mu \zeta^a \tensor{e}{_c^\mu} \wedge \theta^c + \omega^a_{cb} \zeta^b \wedge \theta^c = \left[ e_c(\zeta^a) + \omega^a_{cb} \zeta^b \right] \wedge \theta^c \equiv \tensor{\zeta}{^a_{;c}} \wedge \theta^c
\end{equation}

The same construction for the exterior covariant derivative can be generalized to tensor-valued $ p $-forms as:
\begin{equation}
  \dcov \tensor{\zeta}{^{a_1 \dots a_r}_{b_1 \dots b_s}} = \dd \tensor{\zeta}{^{a_1 \dots a_r}_{b_1 \dots b_s}} + \sum_{j = 1}^r \tensor{\omega}{^{a_j}_c} \wedge \tensor{\zeta}{^{a_1 \dots c \dots a_r}_{b_1 \dots b_s}} - \sum_{j = 1}^s \tensor{\omega}{^c_{b_j}} \wedge \tensor{\zeta}{^{a_1 \dots a_r}_{b_1 \dots c \dots b_s}}
  \label{eq:dcov}
\end{equation}
It is then trivial to show, using \lref{lemma:form-lor-inv}, that a local Lorentz transformation acts as:
\begin{equation}
  \tensor{\zeta}{^{a_1 \dots a_r}_{b_1 \dots b_s}} \mapsto \tensor{\Lambda}{^{a_1}_{c_1}} \dots \tensor{\Lambda}{^{a_r}_{c_r}} \tensor{\zeta}{^{c_1 \dots c_r}_{d_1 \dots d_s}} \tensor{(\Lambda^{-1})}{^{d_1}_{b_1}} \dots \tensor{(\Lambda^{-1})}{^{d_s}_{b_s}}
\end{equation}
\begin{equation}
  \dcov\tensor{\zeta}{^{a_1 \dots a_r}_{b_1 \dots b_s}} \mapsto \tensor{\Lambda}{^{a_1}_{c_1}} \dots \tensor{\Lambda}{^{a_r}_{c_r}} \dcov\tensor{\zeta}{^{c_1 \dots c_r}_{d_1 \dots d_s}} \tensor{(\Lambda^{-1})}{^{d_1}_{b_1}} \dots \tensor{(\Lambda^{-1})}{^{d_s}_{b_s}}
\end{equation}
where the transformation rule $ \zeta_a \mapsto \zeta_b \tensor{(\Lambda^{-1})}{^b_a} $ follows from the requirement of $ \zeta^a \xi_a $ being Lorentz-invariant (as for usual Lorentz transformation).

\subsubsection{Levi--Civita spin connection}

If the connection is torsionless, the first structure equation reduces to:
\begin{equation}
  \dcov \theta^a = 0
  \label{eq:first-cartan-stru-eq}
\end{equation}
For a Levi--Civita connection, a stronger result holds.

\begin{proposition}{}{}
  On a metric manifold with a Levi--Civita connection:
  \begin{equation}
    \omega_{ab} = -\omega_{ba}
    \label{eq:connection-form-symm}
  \end{equation}
\end{proposition}

\begin{proofbox}
  \begin{proof}
    Consider the vielbein condition \eref{eq:vielbeins-def}. For a Levi--Civita connection, compatibility with the metric implies:
    \begin{equation*}
      \begin{split}
        0
        & = \na_X (g_{\mu \nu} \tensor{e}{_a^\mu} \tensor{e}{_b^\nu}) = X^c g_{\mu \nu} \left( \na_{e_c} \tensor{e}{_a^\mu} \tensor{e}{_b^\nu} + \tensor{e}{_a^\mu} \na_{e_c} \tensor{e}{_b^\nu} \right) \\
        & = X^c g_{\mu \nu} \left( \omega^d_{ca} \tensor{e}{_d^\mu} \tensor{e}{_b^\nu} + \omega^d_{cb} \tensor{e}{_a^\mu} \tensor{e}{_d^\nu} \right) = X^c \left( \omega^d_{ca} \eta_{db} + \omega^d_{cb} \eta_{ad} \right)
      \end{split}
    \end{equation*}
    Note that $ \tensor{\omega}{^a_b}(X) = \omega^a_{cb} \theta^c(X) = \omega^a_{cb} X^c $, so that:
    \begin{equation*}
      0 = \tensor{\omega}{^d_a}(X) \eta_{db} + \tensor{\omega}{^d_b}(X) \eta_{ad} = \omega_{ba}(X) + \omega_{ab}(X)
    \end{equation*}
    By the arbitrariness of $ X $, it follows the thesis.
  \end{proof}
\end{proofbox}

\eeref{eq:first-cartan-stru-eq}{eq:connection-form-symm} allow to quickly compute the spin connection of a Levi--Civita connection, as they uniquely define it. Indeed, by antisymmetry $ \omega_{ab} $ has $ \frac{1}{2} n (n-1) $ independent 1-forms, i.e. $ \frac{1}{2} n^2 (n-1) $ independent components. On the other hand, the first structure equation constraints a vector-valued 2-form, which has $ n $ 2-forms with $ \frac{1}{2} n (n-1) $ independent components each, thus posing $ \frac{1}{2} n^2 (n-1) $ constraints: this is the same number as the independent components of $ \omega_{ab} $, therefore the spin connection is uniquely determined.

The Levi--Civita connection is the unique connection such that \eeref{eq:first-cartan-stru-eq}{eq:connection-form-symm} both hold. Moreover, for a Levi--Civita connection $ \dcov (\zeta^a \xi_a) = \dd (\zeta^a \xi_a) $ and $ \dcov \eta_{ab} = 0 $, since $ \tensor{\omega}{_a^b} = - \tensor{\omega}{^b_a} $.

\subsection{Curvature 2-form}

As for the torsion tensor, it is possible to compute the curvature tensor on a vielbein basis as $ \tensor{R}{^a_{bcd}} = \riem(\theta^a , e_b , e_c , e_d) $, which retains the symmetries \pref{prop:riemann-symmetries}.

\begin{definition}{Curvature 2-form}{}
  On a metric manifold with vielbeins, the \bcdef{curvature 2-form} is defined as:
  \begin{equation}
    \tensor{\Omega}{^a_b} \defeq \frac{1}{2} \tensor{R}{^a_{bcd}} \theta^c \wedge \theta^d
    \label{eq:curvature-2-def}
  \end{equation}
\end{definition}

The curvature 2-form is a tensor-valued 2-form, hence $ \tensor{\Omega}{^a_b} $ has $ n^2 $ 2-forms, i.e. $ \frac{1}{2} n^3 (n-1) $ independent components (the same number as the curvature tensor, as expected). It is possible to derive a structure equation for the curvature 2-form too.

\begin{theorem}[before upper = {\tcbtitle}]{Second Cartan structure equation}{}
  \begin{equation}
    \tensor{\Omega}{^a_b} = \dd \tensor{\omega}{^a_b} + \tensor{\omega}{^a_c} \wedge \tensor{\omega}{^c_b}
    \label{eq:second-cartan-stru-eq}
  \end{equation}
\end{theorem}

\begin{proofbox}
  \begin{proof}
    Recalling \eref{eq:riemann-operator-def}, it is possible to write:
    \begin{equation*}
      \tensor{R}{^a_{bcd}} e_a = \riem(e_c , e_d) e_b = \na_{e_c} \na_{e_d} e_b - \na_{e_d} \na_{e_c} e_b - \na_{[e_c , e_d]} e_b
    \end{equation*}
    In general, by linearity of the connection, $ \na_X e_b = \tensor{\omega}{^c_b}(X) e_c $, hence (using $ \na_X f = X(f) $):
    \begin{equation*}
      \begin{split}
        \tensor{R}{^a_{bcd}} e_a
        & = \na_{e_c} [\omega^a_{db} e_a] - \na_{e_d} [\omega^a_{cb} e_a] - \tensor{\omega}{^a_b}([e_c , e_d]) e_a \\
        & = e_c(\omega^a_{db}) e_a + \omega^f_{db} \omega^a_{cf} e_a - e_d(\omega^a_{cb}) e_a - \omega^f_{cb} \omega^a_{df} e_a - \tensor{\omega}{^a_b}([e_c , e_d]) e_a \\
      \end{split}
    \end{equation*}

    \begin{lemma}[before upper = {\tcbtitle}, borderline = {1pt}{0pt}{blue!35!black}]{}{connection-ext-der}
      \begin{equation}
        \dd \tensor{\omega}{^m_n}(e_a , e_b) = e_a(\omega^m_{bn}) - e_b(\omega^m_{an}) - \tensor{\omega}{^m_n}([e_a , e_b])
      \end{equation}
    \end{lemma}

    \begin{proofbox}
      \begin{proof}
        Writing $ \pa_\mu = \tensor{e}{^a_\mu} e_a $ and $ \dd x^\mu = \tensor{e}{_a^\mu} \theta^a $, by linearity of the connection:
        \begin{equation*}
          \dd \tensor{\omega}{^m_n} = \na_{e_a} (\omega^m_{cn} \tensor{e}{^c_\mu}) \tensor{e}{_b^\mu} \theta^a \wedge \theta^b = \frac{1}{2} \left[ \na_{e_a} (\omega^m_{cn} \tensor{e}{^c_\mu}) \tensor{e}{_b^\mu} - \na_{e_b} (\omega^m_{cn} \tensor{e}{^c_\mu}) \tensor{e}{_a^\mu} \right] \theta^a \wedge \theta^b
        \end{equation*}
        Since $ \omega^m_{cn} \tensor{e}{^c_\mu} $ is a scalar function, using the Leibniz rule:
        \begin{equation*}
          \dd \tensor{\omega}{^m_n}(e_a , e_b) = e_a(\omega^m_bn) - e_b(\omega^m_an) + \omega^m_{cn} \left[ e_a(\tensor{e}{^c_\mu}) \tensor{e}{_b^\mu} - e_b(\tensor{e}{^c_\mu}) \tensor{e}{_a^\mu} \right]
        \end{equation*}
        Now, consider the action of the commutator of vielbeins on $ \cm $:
        \begin{equation*}
          \begin{split}
            [e_a , e_b](f)
            & \defeq e_a(e_b(f)) - e_b(e_a(f)) = \tensor{e}{_a^\mu} \pa_\mu (\tensor{e}{_b^\nu} \pa_\nu f) - \tensor{e}{_b^\nu} \pa_\nu (\tensor{e}{_a^\mu} \pa_\mu f) \\
            & = \tensor{e}{_a^\mu} \pa_\mu \tensor{e}{_b^\nu} \pa_\nu f - \tensor{e}{_b^\nu} \pa_\nu \tensor{e}{_a^\mu} \pa_\mu f = \left[ e_a(\tensor{e}{_b^\mu}) \tensor{e}{^c_\mu} - e_b(\tensor{e}{_a^\mu}) \tensor{e}{^c_\mu} \right] e_c(f) \\
          \end{split}
        \end{equation*}
        As $ \tensor{e}{_b^\mu} \tensor{e}{^c_\mu} = \delta^c_b $, by the Leibniz rule $ e_a(\tensor{e}{_b^\mu}) \tensor{e}{^c_\mu} + \tensor{e}{_b^\mu} e_a(\tensor{e}{^c_\mu}) = 0 $, hence:
        \begin{equation}
          ([e_a , e_b])^c = \tensor{e}{_a^\mu} e_b(\tensor{e}{^c_\mu}) - \tensor{e}{_b^\mu} e_a(\tensor{e}{^c_\mu})
        \end{equation}
        which, inserted in the expression for $ \dd \tensor{\omega}{^m_n}(e_a , e_b) $, yields the thesis.
      \end{proof}
    \end{proofbox}

    Using \lref{lemma:connection-ext-der}, the above expression becomes:
    \begin{equation*}
      \tensor{R}{^a_{bcd}} = \dd \tensor{\omega}{^a_b}(e_c , e_d) + \omega^f_{db} \omega^a_{cf} - \omega^f_{cb} \omega^a_{df}
    \end{equation*}
    But $ \tensor{R}{^a_{bcd}} = \tensor{\Omega}{^a_b}(e_c , e_d) $, while:
    \begin{equation*}
      \tensor{\omega}{^a_f} \wedge \tensor{\omega}{^f_b} = \omega^a_{cf} \omega^f_{db} \theta^c \wedge \theta^d = \frac{1}{2} \left[ \omega^f_{db} \omega^a_{cf} - \omega^f_{cb} \omega^a_{df} \right] \theta^c \wedge \theta^d
    \end{equation*}
    thus the proof is complete.
  \end{proof}
\end{proofbox}

Combining the two structure equations, it is possible to find simple expressions for the exterior covariant derivative of the torsion and curvature 2-forms.

\begin{proposition}[before upper = {\tcbtitle}]{}{}
  \begin{equation}
    \dcov T^a = \tensor{\Omega}{^a_b} \wedge e^b
    \qquad \qquad
    \dcov \tensor{\Omega}{^a_b} = 0
    \label{eq:cov-ext-der-torsion-curvature}
  \end{equation}
\end{proposition}

\begin{proofbox}
  \begin{proof}
    From \eref{eq:first-cartan-stru-eq}, taking the exterior derivative and using \eref{eq:ext-der-prod-rule}:
    \begin{equation*}
      \dd T^a = \dd \tensor{\omega}{^a_b} \wedge \theta^b - \tensor{\omega}{^a_b} \wedge \dd \theta^b = \dd \tensor{\omega}{^a_b} \wedge \theta^b - \tensor{\omega}{^a_b} \wedge (T^b - \tensor{\omega}{^b_c} \wedge \theta^c)
    \end{equation*}
    Rearranging the terms and inserting \eref{eq:second-cartan-stru-eq}:
    \begin{equation*}
      \dd T^a + \tensor{\omega}{^a_b} \wedge T^a = \left( \dd \tensor{\omega}{^a_b} + \tensor{\omega}{^a_c} \wedge \tensor{\omega}{^c_b} \right) \wedge \theta^b
      \quad \implies \quad
      \dcov T^a = \tensor{\Omega}{^a_b} \wedge e^b
    \end{equation*}
    Now, take the exterior derivative of \eref{eq:second-cartan-stru-eq} and apply \eref{eq:ext-der-prod-rule}:
    \begin{equation*}
      \dd \tensor{\Omega}{^a_b} = \dd \tensor{\omega}{^a_c} \wedge \tensor{\omega}{^c_b} - \tensor{\omega}{^a_c} \wedge \dd \tensor{\omega}{^c_b} = \left( \tensor{\Omega}{^a_c} - \tensor{\omega}{^a_d} \wedge \tensor{\omega}{^d_c} \right) \wedge \tensor{\omega}{^c_b} - \tensor{\omega}{^a_c} \wedge \left( \tensor{\Omega}{^c_b} - \tensor{\omega}{^c_d} \wedge \tensor{\omega}{^d_b} \right)
    \end{equation*}
    The second and fourth terms cancel by renaming the mute indices, thus:
    \begin{equation*}
      0 = \dd \tensor{\Omega}{^a_b} + \tensor{\omega}{^a_c} \wedge \tensor{\Omega}{^c_b} - \tensor{\Omega}{^a_c} \wedge \tensor{\omega}{^c_b} = \dcov \tensor{\Omega}{^a_b}
    \end{equation*}
    since $ \tensor{\Omega}{^a_c} \wedge \tensor{\omega}{^c_b} = \tensor{\omega}{^c_b} \wedge \tensor{\Omega}{^a_c} $ by \eref{eq:form-skew-symmetry}.
  \end{proof}
\end{proofbox}

\subsubsection{Bianchi identities}

The Bianchi identities are useful identities which hold for a torsionless connection and put symmetry constraints on the curvature tensor.

\begin{theorem}{Bianchi identities}{}
  For a torsionless connection:
  \begin{equation}
    \tensor{R}{^a_{[bcd]}} = 0
    \qquad \qquad
    \tensor{R}{^a_{b[cd;k]}} = 0
    \label{eq:bianchi-vielbein}
  \end{equation}
\end{theorem}

\begin{proofbox}
  \begin{proof}
    Setting $ T^a \equiv 0 $ in the first \eref{eq:cov-ext-der-torsion-curvature}:
    \begin{equation*}
      0 = \tensor{\Omega}{^a_b} \wedge \theta^b = \frac{1}{2} \tensor{R}{^a_{bcd}} \theta^c \wedge \theta^d \wedge \theta^b = \frac{1}{2} \tensor{R}{^a_{[bcd]}} \theta^b \wedge \theta^c \wedge \theta^d
    \end{equation*}
    which is the first Bianchi identity. Then, from the second \eref{eq:cov-ext-der-torsion-curvature}:
    \begin{equation*}
      0 = \dcov \tensor{\Omega}{^a_b} = \frac{1}{2} \dcov \left[ \tensor{R}{^a_{bcd}} \theta^c \wedge \theta^d \right] = \frac{1}{2} \left[ \tensor{R}{^a_{bcd;k}} \theta^k \wedge \theta^c \wedge \theta^d + \tensor{R}{^a_{bcd}} \left( \dcov \theta^c \wedge \theta^d - \theta^c \wedge \dcov \theta^d \right) \right]
    \end{equation*}
    But $ \dcov \theta^a = T^a \equiv 0 $, hence this equation is equivalent to the second Bianchi identity.
  \end{proof}
\end{proofbox}

In the vielbein basis, the Ricci tensor is defined as the contraction $ R_{bd} \equiv \tensor{R}{^a_{bad}} $, while the Ricci scalar is the trace $ R \equiv \eta^{ab} R_{ab} $. In the same way, the Einstein tensor is defined as $ G_{ab} \equiv R_{ab} - \frac{1}{2} \eta_{ab} R $.

\begin{corollary}{Einstein--Bianchi identity}{}
  For a Levi--Civita connection:
  \begin{equation}
    \tensor{G}{_{ab}^{;b}} = 0
  \end{equation}
\end{corollary}

\begin{proofbox}
  \begin{proof}
    From the second Bianchi identity \eref{eq:bianchi-vielbein}, using the antisymmetry $ \tensor{R}{^a_{bcd}} = \tensor{R}{^a_{b[cd]}} $:
    \begin{equation*}
      R_{bd;k} + \tensor{R}{^a_{bka;d}} + \tensor{R}{^a_{bdk;a}} = 0
      \quad \implies \quad
      R_{bd;k} - R_{bk;d} + \tensor{R}{^a_{bdk;a}} = 0
    \end{equation*}
    Now, contract with $ \eta^{bd} $:
    \begin{equation*}
      R_{;k} - \tensor{R}{^b_{k;b}} + \tensor{R}{^{ad}_{dk;a}} = 0
      \quad \implies \quad
    \end{equation*}

    \begin{lemma}[borderline = {1pt}{0pt}{blue!35!black}]{Symmetries of Riemann tensor}{riemann-symmetry-vielbein}
      For a Levi--Civita connection:
      \begin{equation}
        R_{abcd} = R_{[ab][cd]}
      \end{equation}
    \end{lemma}

    \begin{proofbox}
      \begin{proof}
        By \eref{eq:curvature-2-def} $ \tensor{R}{^a_{bcd}} = \tensor{\Omega}{^a_b}(e_c , e_d) $, thus to compute this expression first consider:
        \begin{equation*}
          \dd \tensor{\omega}{^a_b} = \dd \omega^a_{mb} \wedge \theta^m + \omega^a_{mb} \dd \theta^m = \dd \omega^a_{mb} \wedge \theta^m - \omega^a_{mb} \tensor{\omega}{^m_k} \wedge \theta^k
        \end{equation*}
        where \eref{eq:first-cartan-stru-eq} was used. Then, since $ \dd f = \pa_\mu f \tensor{e}{_a^\mu} \theta^a = e_a(f) \theta^a \,\, \forall f \in \cm $:
        \begin{equation*}
          \dd \tensor{\omega}{^a_b} = e_k(\omega^a_{mb}) \theta^k \wedge \theta^m - \omega^a_{mb} \tensor{\omega}{^m_k} \wedge \theta^k = \left[ e_m(\omega^a_{kb}) - \omega^a_{cb} \omega^c_{mk} \right] \theta^m \wedge \theta^k
        \end{equation*}
        Therefore, form \eref{eq:second-cartan-stru-eq}:
        \begin{equation*}
          \tensor{\Omega}{^a_b} = \left[ e_m(\omega^a_{kb}) - \omega^a_{cb} \omega^c_{mk} + \omega^a_{mc} \omega^c_{kb} \right] \theta^m \wedge \theta^k
        \end{equation*}
        Now, to evaluate $ \tensor{\Omega}{^a_b}(e_c , e_d) $, use:
        \begin{equation*}
          \theta^m \wedge \theta^k (e_c , e_d) \equiv
          \begin{vmatrix}
            \theta^m(e_c) & \theta^m(e_d) \\ \theta^k(e_c) & \theta^k(e_d)
          \end{vmatrix}
          = \delta^m_c \delta^k_d - \delta^m_d \delta^k_c \eqdef \delta^{mk}_{cd}
        \end{equation*}
        Hence, defining $ \omega_{abc} \equiv \eta_{ak} \omega^k_{bc} $, the curvature tensor is found to be:
        \begin{equation*}
          R_{abcd} = \left[ e_m(\omega_{akb}) - \omega_{agb} \omega^g_{mk} + \omega_{amg} \omega^g_{kb} \right] \delta^{mk}_{cd}
        \end{equation*}
        Then, since $ \delta^{mk}_{cd} = \delta^{[mk]}_{[cd]} $ and, recalling \eref{eq:connection-form-symm} for a Levi--Civita connection, $ \omega_{ab} = - \omega_{ba} \implies \omega_{acb} = \omega_{[a|c|b]} $, this is equivalente to the thesis.
      \end{proof}
    \end{proofbox}

    Using \lref{lemma:riemann-symmetry-vielbein}, the above expression reduces to:
    \begin{equation*}
      R_{;k} = 2 \tensor{R}{^a_{k;a}}
    \end{equation*}
    Then, the thesis is trivial from the metric-compatibility of the connection.
  \end{proof}
\end{proofbox}

Note that \lref{lemma:riemann-symmetry-vielbein} allows to explicitly compute the curvature tensor from connection coefficients:
\begin{equation}
  R_{abcd} = e_c(\omega_{adb}) - e_d(\omega_{acb}) - \omega_{agb} \left( \omega^g_{cd} - \omega^g_{dc} \right) + \omega_{acg} \omega^g_{db} - \omega_{adg} \omega^g_{cb}
\end{equation}
Moreover, by multilinearity, the curvature tensor in the vielbein basis retains the symmetry $ R_{abcd} = R_{cdab} $. This, together with $ R_{abcd} = R_{[ab][cd]} $, means that the Riemann tensor can be thought as having two symmetrized multi-indices $ [ab] $ and $ [cd] $, hence the number of independent components is $ \frac{1}{2} N (N+1) $, where $ N $ is the number of possible values for a single multi-index, i.e. $ N = \frac{1}{2} n (n-1) $. However, one has to account for the constraint posed by the first Bianchi identity: since there is one such constraint for each choise of four distinc indices, the total number of independent components of the Riemann tensor is:
\begin{equation*}
  f_\text{Riemann} = \frac{1}{2} \frac{n (n-1)}{2} \left( \frac{n(n-1)}{2} + 1 \right) - \binom{n}{4} = \frac{1}{8} (n-1) (n (n-1) + 2) - \frac{1}{24} n (n-1) (n-2) (n-3)
\end{equation*}
Simplifying this expression yields the classical result:
\begin{equation}
  f_\text{Riemann} = \frac{n^2 (n^2-1)}{12}
\end{equation}
In particular, on $ 4 $-dimensional manifolds this analysis yields the separation of the $ 20 $ degrees of freedom of the Riemann tensor into the $ 10 $ degrees of freedom of the Ricci tensor (``symmetric" part) and the $ 10 $ of the Weyl tensor (``antisymmetric" part).

Note that, if the curvature tensor is computed on a non-coordinate and non-tetrad basis, then it is not possible to impose the symmetry $ R_{abcd} = R_{cdab} $ and the unconstrained components rise to $ 36 $ (i.e. $ N^2 $): the spurious $ 16 $ degrees of freedom thus need to be eliminated with additional constraints, depending on the formalism employed.










